\begin{figure}[t]
\centering
\resizebox{\linewidth}{!}{%
\begin{tikzpicture}[
  font=\small,
  box/.style={draw=black!45, rounded corners=2pt, line width=0.55pt, align=center, minimum width=2.85cm, minimum height=0.82cm, fill=white},
  source/.style={box, fill=blue!5},
  evidence/.style={box, fill=cyan!5},
  knowledge/.style={box, fill=orange!9},
  model/.style={box, fill=green!8},
  output/.style={box, fill=purple!6},
  lane/.style={draw=black!18, rounded corners=4pt, fill=black!2},
  arrow/.style={-{Latex[length=2.1mm]}, line width=0.65pt, draw=black!70},
  softarrow/.style={-{Latex[length=1.8mm]}, line width=0.55pt, draw=black!45}
]
\node[source] (q) at (0,1.45) {\textbf{输入}\\医学问题 $q$};
\node[evidence] (retrieval) at (3.35,1.45) {\textbf{候选生成}\\BM25、稠密、fielded BM25、\\MedCPT 双编码器};
\node[evidence] (pool) at (6.95,1.45) {\textbf{Top-$M$ 证据池}\\验证集选择的\\加权 RRF};
\node[knowledge] (kg) at (3.35,-0.15) {\textbf{医学约束}\\实体、MeSH 层级、\\PrimeKG 辅助链接};
\node[knowledge] (hg) at (6.95,-0.15) {\textbf{局部超图}\\问题、段落、文档、\\实体与概念节点};
\node[model] (features) at (10.45,-0.15) {\textbf{可解释特征}\\语义分数、层级、\\扩散、中心性};
\node[model] (ltr) at (10.45,1.45) {\textbf{查询分组 LTR}\\LambdaMART 打分};
\node[output] (out) at (13.75,1.45) {\textbf{输出}\\排序证据 top-$k$\\与可选路径};

\begin{scope}[on background layer]
  \node[lane, fit=(q)(retrieval)(pool), inner sep=8pt] (laneA) {};
  \node[lane, fit=(kg)(hg)(features), inner sep=8pt] (laneB) {};
\end{scope}
\node[anchor=west, font=\footnotesize\bfseries, text=black!62] at ($(laneA.north west)+(0.12,-0.18)$) {检索层};
\node[anchor=west, font=\footnotesize\bfseries, text=black!62] at ($(laneB.north west)+(0.12,-0.18)$) {证据控制层};

\draw[arrow] (q) -- (retrieval);
\draw[arrow] (retrieval) -- (pool);
\draw[arrow] (pool) -- (ltr);
\draw[arrow] (ltr) -- (out);
\draw[softarrow] (kg) -- (hg);
\draw[softarrow] (pool) -- (hg);
\draw[softarrow] (q.south) |- (hg.west);
\draw[softarrow] (hg) -- (features);
\draw[softarrow] (features) -- (ltr);
\end{tikzpicture}}
\caption{KCH-MedRank 的论文风格方法总览。候选生成阶段提供高召回证据池；证据控制层将生物医学约束和局部超图结构转换为可解释特征；查询分组排序器输出用于下游 medical RAG 的最终证据顺序。}
\label{fig:method-overview}
\end{figure}
